\cleardoublepage
\phantomsection
\tableofcontents

\clearpage
\phantomsection
\addcontentsline{toc}{chapter}{Índice de figuras}
\listoffigures

\clearpage
\phantomsection
\addcontentsline{toc}{chapter}{Índice de tablas}
\listoftables

\clearpage
\phantomsection
\addcontentsline{toc}{chapter}{Índice de algoritmos}
\listofalgorithms

\clearpage
\phantomsection
\addcontentsline{toc}{chapter}{Índice de ecuaciones}
\listofequations

\clearpage
\phantomsection
\addcontentsline{toc}{chapter}{Índice de anexos}
\listofanexos

% ---------------- Lista de abreviaturas y siglas ----------------
\clearpage
\chapter*{Índice de abreviaturas}
\addcontentsline{toc}{chapter}{Índice de abreviaturas}
\begingroup\singlespacing\small
\begin{longtable}{@{}p{3.3cm}p{11.5cm}@{}}
ARI & \textit{Adjusted Rand Index} (índice de Rand ajustado). \\
CRISP-DM & \textit{Cross-Industry Standard Process for Data Mining}. \\
CSR & \textit{Complete Spatial Randomness} (aleatoriedad espacial completa). \\
DBSCAN & \textit{Density-Based Spatial Clustering of Applications with Noise}. \\
ENAPRES & Encuesta Nacional de Programas Presupuestales. \\
GIS & \textit{Geographic Information System} (equivalente en inglés de SIG). \\
HDBSCAN & \textit{Hierarchical Density-Based Spatial Clustering of Applications with Noise}. \\
INEI & Instituto Nacional de Estadística e Informática. \\
KDE & \textit{Kernel Density Estimation} (estimación de densidad kernel). \\
PNP & Policía Nacional del Perú. \\
SIG & Sistema de Información Geográfica. \\
SIPCOP-M & Sistema de registro de las ocurrencias de seguridad atendidas por el serenazgo de la Municipalidad Distrital de San Sebastián. \\
UNSAAC & Universidad Nacional de San Antonio Abad del Cusco. \\
UTM & \textit{Universal Transverse Mercator}. \\
WGS84 & \textit{World Geodetic System 1984}. \\
\end{longtable}
\endgroup

% ---------------- Lista de símbolos ----------------
\clearpage
\chapter*{Índice de símbolos}
\addcontentsline{toc}{chapter}{Índice de símbolos}
\begingroup\singlespacing\small
\begin{longtable}{@{}p{3.3cm}p{11.5cm}@{}}
$\varepsilon$ & Radio de vecindad de DBSCAN (distancia máxima entre dos puntos vecinos). \\
$\mathrm{MinPts}$ & Número mínimo de puntos dentro del radio $\varepsilon$ para considerar un punto como núcleo. \\
$N_{\varepsilon}(p)$ & Conjunto de puntos a distancia menor o igual que $\varepsilon$ del punto $p$. \\
$d(p,q)$ & Distancia entre los puntos $p$ y $q$. \\
$R$ & Radio medio de la Tierra (aprox. 6371 km). \\
$\varphi,\ \lambda$ & Latitud y longitud, respectivamente, expresadas en radianes. \\
$h$ & Ancho de banda (\textit{bandwidth}) del estimador de densidad kernel. \\
$K(\cdot)$ & Función kernel. \\
$\hat{f}(x)$ & Densidad estimada de incidentes en la ubicación $x$. \\
$\chi^{2}$ & Estadístico chi-cuadrado de Pearson. \\
$O_i,\ E_i$ & Frecuencia observada y esperada de la categoría $i$. \\
$H_0,\ H_1$ & Hipótesis nula y alterna. \\
$\alpha$ & Nivel de significancia. \\
$s(i)$ & Coeficiente de silueta del punto $i$. \\
$a(i),\ b(i)$ & Distancia media del punto $i$ a los de su grupo y al grupo vecino más cercano. \\
$V$ & Coeficiente V de Cramér (tamaño del efecto). \\
$B$ & Número de simulaciones Monte Carlo. \\
$m_j$ & Número de ocurrencias del grupo espacial $j$. \\
$M_b$ & Tamaño del mayor grupo en la simulación $b$. \\
$p_j$ & Valor $p$ Monte Carlo del grupo $j$. \\
\end{longtable}
\endgroup

% ---------------- Glosario ----------------
\clearpage
\chapter*{Glosario de términos}
\addcontentsline{toc}{chapter}{Glosario de términos}
\begingroup\singlespacing\small
\begin{longtable}{@{}p{3.3cm}p{11.5cm}@{}}
\textbf{Agrupamiento espacial} & Técnica de análisis que reúne eventos geográficamente cercanos en grupos (\textit{clusters}) según un criterio de proximidad o densidad. \\
\textbf{Densidad kernel} & Método que estima una superficie continua de intensidad de eventos a partir de puntos, ponderando cada evento según su distancia. \\
\textbf{Epsilon ($\varepsilon$)} & Parámetro de DBSCAN que define el radio de búsqueda de vecinos alrededor de un punto. \\
\textbf{Georreferenciación} & Asignación de coordenadas geográficas a un registro para ubicarlo en un sistema de referencia espacial. \\
\textbf{Haversine} & Fórmula que calcula la distancia sobre una esfera entre dos puntos dados por latitud y longitud. \\
\textbf{Grupo espacial (clúster)} & Conjunto de ocurrencias que un algoritmo de agrupamiento reúne por su proximidad o densidad. \\
\textbf{Hotspot} & Grupo espacial cuyo número de ocurrencias supera lo esperado bajo aleatoriedad espacial completa, según un contraste estadístico. \\
\textbf{Ocurrencia de seguridad ciudadana} & Hecho de seguridad atendido y registrado por el serenazgo; incluye hechos que no son necesariamente delitos. \\
\textbf{Mapa de calor} & Representación cartográfica que usa gradientes de color para mostrar la intensidad de un fenómeno. \\
\textbf{MinPts} & Parámetro de DBSCAN que fija el mínimo de puntos requeridos en la vecindad para formar un grupo. \\
\textbf{Patrón espacio-temporal} & Regularidad en la distribución geográfica y en el momento de ocurrencia (turno, día, mes) de los eventos. \\
\textbf{Ruido} & Puntos que DBSCAN no asigna a ningún grupo por estar en zonas de baja densidad. \\
\textbf{Serenazgo} & Servicio municipal de seguridad ciudadana que realiza patrullaje y prevención en el distrito. \\
\textbf{Turno} & Franja horaria de operación del serenazgo (por ejemplo, mañana, tarde o noche). \\
\end{longtable}
\endgroup

% Los longtable de los índices incrementan el contador de tablas; se reinicia.
\setcounter{table}{0}
\clearpage
